\subsection{Two angular bounds do not give a better pairwise strip power}
\label{angularobs:section}

The positive selection in Section~\ref{angselect:section} leaves a
far-source collision estimate to prove. The next example shows why that
estimate must use the distribution of source pairs, rather than just the
two selected angular marginals at an individual pair.

\begin{proposition}\label{angularobs:proposition}
For every \(1<s<2\), there are separated compactly supported
\(s\)-Ahlfors regular probabilities \(\mu,\nu\) such that the radial
projection of \(\nu\) from every point of \(\supp\mu\) has an
angular density bounded by one fixed constant \(A_0\). Nevertheless,
there are source points \(x_+,x_-\), at distance two, for which
\[
 \nu\{y:\bigl||x_+-y|-|x_--y|\bigr|\le\varepsilon\}
 \asymp\varepsilon^{s-1}.
\]
For every sufficiently small \(\varepsilon\), the same lower bound
holds on a set of far source pairs of positive \(\mu\times\mu\)
mass. Thus no uniform pairwise bound with a larger power of
\(\varepsilon\) follows from these hypotheses, even when both angular
selection factors are retained.
\end{proposition}

\begin{proof}
Let \(\kappa\) be an \((s-1)\)-Ahlfors regular Cantor probability
on a subset of \([0,1/10]\) containing zero, and let \(\lambda\)
be an \((s/2)\)-Ahlfors regular probability with the same support
conventions. Two-map self-similar Cantor measures with ratios
\(2^{-1/(s-1)}\) and \(2^{-2/s}\), followed by fixed dilations,
give such measures. Put \(x_+=(1,0)\), \(x_-=(-1,0)\), and set
\[
 \mu=\tfrac12\bigl((x_++\cdot)_*(\lambda\times\lambda)
                 +(x_-+\cdot)_*(\lambda\times\lambda)\bigr),
 \qquad
 \nu=\kappa\times\mathcal L^1|_{[2,3]}.
\]
Products are \(s\)-Ahlfors regular by comparing balls with product
intervals of comparable radius. The separated union forming \(\mu\)
preserves this property. Thus both supports have Hausdorff and packing
dimension \(s\), and the source and pin supports are separated.

For a source \(x\), fix a horizontal pin coordinate \(t\in[0,1/10]\).
The angle of \((t,v)-x\), as \(v\) ranges over \([2,3]\), is
injective and has derivative in absolute value
\[
 \frac{|t-x_1|}{(t-x_1)^2+(v-x_2)^2}.
\]
This is bounded above and below by positive constants uniformly in all
the specified coordinates. The angular pushforward of vertical Lebesgue
measure therefore has a uniformly bounded density. Integrating against
\(\kappa(t)\) gives jointly measurable angular densities with the
same bound \(A_0\). With these representatives, any fixed threshold
\(A\ge A_0\) retains every original pair in the angular selection.

For the two distinguished source points and \(y=(t,v)\), the exact
squared-distance identity gives
\[
 \bigl||x_+-y|-|x_--y|\bigr|
 =\frac{4|t|}{|x_+-y|+|x_--y|}.
\]
The denominator stays between positive fixed constants. Upper and lower
regularity of \(\kappa\) at zero prove the claimed power.

This failure is not confined to one null source pair at a fixed scale.
If \(x\in B(x_+,c_0\varepsilon)\cap\supp\mu\) and
\(x'\in B(x_-,c_0\varepsilon)\cap\supp\mu\), the distance
difference changes by at most \(2c_0\varepsilon\). With \(c_0\)
small, all pins with \(t\le c\varepsilon\) still collide.
Their mass is at least \(c'\varepsilon^{s-1}\), and this source-pair
set has product mass at least \(c''\varepsilon^{2s}\).
All its source separations exceed one for small \(\varepsilon\).
A uniform upper bound \(C\varepsilon^\theta\), \(\theta>s-1\),
fails on this positive-mass set once the resolution is small enough.
\end{proof}

The displayed neighborhoods of source pairs contribute only a lower
bound of order \(\varepsilon^{3s-2}\) to the normalized collision
integral. This tends to zero. Consequently the example supplies
\emph{no} counterexample to the source-pair-averaged far criterion,
nor to positive-length distances. Its conclusion is precisely that
bounded angular marginals alone cannot improve the pairwise strip
estimate by a positive power.

There is also no product-density conclusion hidden in the two angular
bounds. On the pin rectangle the map
\(y\mapsto(\arg(y-x_+),\arg(y-x_-))\) has Jacobian magnitude
\[
 \frac{2v}{|y-x_+|^2|y-x_-|^2},\qquad y=(t,v),
\]
bounded below by a positive constant. Its joint angular law is still
singular with respect to two-dimensional Lebesgue measure: the source
of this map is \(\supp\nu\), of dimension \(s<2\), and the map
is Lipschitz there. Both one-dimensional marginals have bounded
densities. Thus even transverse angular coordinates cannot be treated
as independent or as having a bounded joint density.
